\documentclass[11pt,a4paper]{article}
\usepackage[T1]{fontenc}
\usepackage[utf8]{inputenc}
\usepackage{lmodern,microtype}
\usepackage[margin=25mm]{geometry}
\usepackage{amsmath,amssymb,amsthm,booktabs,array}
\usepackage{xcolor,hyperref,fancyhdr,enumitem}
\definecolor{navy}{RGB}{28,49,75}
\hypersetup{colorlinks=true,linkcolor=navy,urlcolor=navy,citecolor=navy,pdftitle={A three-state thermal counterexample to a heat-current kurtosis conjecture},pdfauthor={Dev D. Goyal},pdfsubject={An exact thermal spectral counterexample with a Lean formalization}}
\pagestyle{fancy}\fancyhf{}\fancyhead[L]{\small Thermal heat-current kurtosis}\fancyhead[R]{\small 27 October 2026}\fancyfoot[C]{\thepage}
\setlength{\headheight}{14pt}
\setlength{\parskip}{5pt}\setlength{\parindent}{0pt}
\setlist{nosep,leftmargin=*}
\newtheorem{theorem}{Theorem}
\newtheorem{lemma}{Lemma}
\newcommand{\R}{\mathbb R}\newcommand{\C}{\mathbb C}
\newcommand{\K}{\mathcal K}
\title{\color{navy}A three-state thermal counterexample\protect\\to a heat-current kurtosis conjecture}
\author{Dev D. Goyal\\[3pt]\small\href{https://orcid.org/0009-0009-9132-744X}{ORCID: 0009-0009-9132-744X}}
\date{27 October 2026}
\begin{document}
\maketitle
\begin{abstract}
An explicit three-state, two-bath thermal rate model has a strictly negative fourth-to-second spectral heat-cumulant ratio. Its total off-diagonal rates are positive integers. The energies are distinct, and every supported bath transition satisfies local detailed balance at a positive temperature. The exact ratio is
\[
\K^h=-\frac{3470947982622836}{176627429595846691}<0.
\]
This contradicts the proposed thermal nonnegativity bound, Eq.~(19) of Ptaszy\'nski, Physical Review E \textbf{106}, 024119 (2022), in the finite tilted-generator formulation used in that paper. The proof constructs a $C^4$ local eigenvalue branch, proves that it has the largest real part, and computes its derivatives. Lean~4 verifies these finite spectral statements with no admitted proofs or additional axioms. The formal development does not include the long-time stochastic theorem that identifies this branch with a cumulant generating function.
\end{abstract}

\section{Target, definitions, and scope}
Ptaszy\'nski~\cite{ptaszynski2022} proposes $\K^h=C_4^h/C_2^h\geq0$ for classical Markov systems driven only by temperature differences. The relevant model has fixed state energies, bath-resolved transition rates obeying local detailed balance, and an exponential heat-counting tilt. Its steady-state cumulants are computed as derivatives at zero of the dominant eigenvalue of the tilted generator. This paper uses that spectral formulation.

A tilted generator multiplies each off-diagonal transition rate by an exponential that records the heat transferred on that transition. Its dominant eigenvalue is the eigenvalue with largest real part. The derivatives at zero of this eigenvalue branch are called spectral cumulants.

Here $C_4/C_2$ is the paper's normalized cumulant ratio. It differs from the ordinary fourth central moment divided by the square of the variance. The row-generator convention is used: $q_{ij}$ is the transition rate from state $i$ to state $j$. Ptaszy\'nski uses the transposed, column-generator convention. Transposition preserves the eigenvalues and hence the spectral cumulants.

\textbf{Formal scope.} The verified theorem is about the finite rate matrices, thermal detailed balance, their complex eigenvalues, and real derivatives of a local dominant eigenvalue branch. It does not construct a trajectory probability space or formalize the standard theorem identifying the spectral function with the long-time cumulant generating function of a continuous-time Markov process. It makes no empirical claim about a physical device. This scope is sufficient to contradict the explicitly specified spectral inequality.

\section{Explicit thermal model and result}
Let the state set be $\{0,1,2\}$, with energies and inverse temperatures
\begin{equation}
E=(0,2,1),\qquad \beta_h=\frac{\log2}{2},\qquad \beta_c=5\log2.
\label{eq:temperatures}
\end{equation}
Boltzmann's constant and the energy unit are set to one. The two bath rate matrices below contain only off-diagonal rates:
\begin{equation}
H=\begin{pmatrix}0&10&0\\20&0&0\\0&0&0\end{pmatrix},\qquad
C=\begin{pmatrix}0&0&1\\0&0&32\\32&1&0\end{pmatrix}.
\label{eq:baths}
\end{equation}
Their off-diagonal sum, with escape rates on the diagonal, gives
\begin{equation}
Q=\begin{pmatrix}-11&10&1\\20&-52&32\\32&1&-33\end{pmatrix}.
\label{eq:generator}
\end{equation}
The supported reversible transitions are:
\begin{center}
\begin{tabular}{llll}
\toprule
Bath & Transition & Forward/reverse rates & Forward heat into cold bath\\
\midrule
Hot & $0\leftrightarrow1$ & $10,20$ & $0$\\
Cold & $1\leftrightarrow2$ & $32,1$ & $+1$\\
Cold & $2\leftrightarrow0$ & $32,1$ & $+1$\\
\bottomrule
\end{tabular}
\end{center}
An omitted bath channel is omitted in both directions. Thus each bath couples selectively to particular energy transitions. Such selective couplings are part of established thermal three-level models~\cite{segal2018}.

The cold heat increment on a cold jump $i\to j$ is $E_i-E_j$, and it is zero on a hot jump. The associated tilted matrix is
\begin{equation}
Q_t=\begin{pmatrix}
-11&10&e^{-t}\\
20&-52&32e^t\\
32e^t&e^{-t}&-33
\end{pmatrix}.
\label{eq:tilt}
\end{equation}
The diagonal entries retain their original escape rates.

\begin{theorem}[Thermal spectral counterexample]
The displayed matrices satisfy nonnegative bath rates, strictly positive total off-diagonal rates, zero row sums, and local detailed balance for both positive inverse temperatures in~\eqref{eq:temperatures}. There exists a real function $\chi$, $C^4$ at zero, with $\chi(0)=0$, such that for all sufficiently small real $t$, $\chi(t)$ is an eigenvalue of $Q_t$ and every distinct complex eigenvalue has strictly smaller real part. Its derivatives $C_n=\chi^{(n)}(0)$ satisfy
\begin{align}
C_1&=\frac{2920}{341}, &
C_2&=\frac{447871920}{39651821},\label{eq:c12}\\
C_3&=\frac{43039276005280}{4610753397701}, &
C_4&=-\frac{833027515829480640}{3753001110866489867}.
\label{eq:c34}
\end{align}
Consequently, $C_2>0$, $C_4<0$, and
\begin{equation}
\frac{C_4}{C_2}=-\frac{3470947982622836}{176627429595846691}
\approx-0.019651239847428847.
\label{eq:ratio}
\end{equation}
\end{theorem}

\section{Complete proof}
\subsection{Thermal admissibility}
Every row of $Q$ sums to zero and every off-diagonal entry is positive. Both bath matrices have nonnegative entries. The inverse temperatures satisfy $0<\beta_h<\beta_c$ and $\beta_c=10\beta_h$. The energies are pairwise distinct.

The ratio form of local detailed balance is imposed on supported bath edges; its multiplicative form also includes channels absent in both directions. Written without division by a possibly absent rate, it is
\begin{equation}
k^\alpha_{ij}=k^\alpha_{ji}\exp[-\beta_\alpha(E_j-E_i)].
\label{eq:ldb}
\end{equation}
For the hot excitation $0\to1$, $H_{01}/H_{10}=1/2=e^{-\beta_h 2}$. For either cold relaxation, $1\to2$ or $2\to0$, the energy change is $-1$, and the forward/reverse rate ratio is $32=e^{\beta_c}$. The reverse relations follow by inversion; for unsupported channels both sides of~\eqref{eq:ldb} vanish. Thus~\eqref{eq:ldb} holds for every state pair and both baths.

All energies, temperatures, and rates are fixed. The data involve no chemical potential, external work protocol, negative temperature, or additional nonthermal affinity. The total transition graph has a positive rate between every pair of distinct states.

For an ancillary exact check, the stationary row vector is
\begin{equation}
\pi=\frac{1}{2387}(1684,331,372),\qquad \pi Q=0,\qquad \sum_i\pi_i=1,
\label{eq:stationary}
\end{equation}
and all three entries are positive. These identities are also included in the formal source, though the proof of the counterexample does not require them.

\subsection{Exact characteristic polynomial}
Define
\begin{equation}
P(s)=s^3+96s^2+2387s,\qquad
g(t)=10240(e^{2t}-1)+20(e^{-2t}-1).
\label{eq:pg}
\end{equation}
Direct expansion of the $3\times3$ determinant gives
\begin{equation}
\det(sI-Q_t)=P(s)-g(t).
\label{eq:characteristic}
\end{equation}
One way to check the expansion before substituting exponentials is to replace $e^t,e^{-t}$ by independent variables $z,w$. The determinant then equals
\[
P(s)-10240(z^2-1)-20(w^2-1)+(zw-1)(-64s-2016).
\]
Substitution uses $zw=1$, $z^2=e^{2t}$, and $w^2=e^{-2t}$. The formal proof verifies the determinant over $\C$ as well as over $\R$, so subsequent conclusions address every complex eigenvalue.

\subsection{Existence and spectral identification of the branch}
Since $P(0)=0$ and $P'(0)=2387\ne0$, the $C^4$ inverse-function theorem provides a local inverse $P^{-1}_{\rm loc}$ at zero. Set
\[
\chi(t)=P^{-1}_{\rm loc}(g(t)).
\]
Then $\chi(0)=0$, $\chi$ is $C^4$ at zero, and
\begin{equation}
P(\chi(t))=g(t)
\label{eq:implicit}
\end{equation}
for all sufficiently small $t$. Continuity permits restriction to a neighborhood where $-1<\chi(t)<1$.

The root being differentiated must be identified. Write $\lambda=\chi(t)$. Exact polynomial division gives
\begin{equation}
P(s)-P(\lambda)=(s-\lambda)
\bigl[s^2+(96+\lambda)s+(2387+96\lambda+\lambda^2)\bigr].
\label{eq:factor}
\end{equation}
For $|\lambda|<1$, the discriminant of the quadratic factor is
\[
(96+\lambda)^2-4(2387+96\lambda+\lambda^2)
=-332-192\lambda-3\lambda^2<0.
\]
Its roots are therefore nonreal conjugates with real part $-(96+\lambda)/2<\lambda$. Equation~\eqref{eq:characteristic} shows that $\lambda$ is an eigenvalue of $Q_t$; equation~\eqref{eq:factor} proves strict dominance over all distinct eigenvalues. The Lean proof implements the same argument directly with complex real and imaginary parts and the standard eigenvalue predicate.

\subsection{Four derivatives of the spectral branch}
At zero, the derivatives needed for~\eqref{eq:implicit} are
\[
P'(0)=2387,\qquad P''(0)=192,\qquad P'''(0)=6,
\]
and
\[
g'(0)=20440,\quad g''(0)=41040,\quad
g'''(0)=81760,\quad g^{(4)}(0)=164160.
\]
Applying the chain and product rules successively to~\eqref{eq:implicit} yields
\begin{align}
2387C_1&=20440,\nonumber\\
2387C_2+192C_1^2&=41040,\nonumber\\
2387C_3+576C_1C_2+6C_1^3&=81760,\nonumber\\
2387C_4+768C_1C_3+576C_2^2+36C_1^2C_2&=164160.
\label{eq:recurrence}
\end{align}
The nonzero coefficient $2387$ makes this a triangular system. Exact rational elimination gives~\eqref{eq:c12}--\eqref{eq:c34}; division gives~\eqref{eq:ratio}. The numerators and denominators immediately establish $C_2>0$ and $C_4<0$, completing the proof.

\textbf{What is checked formally in the differentiation step.} The source proves formulas for the first four \texttt{iteratedDeriv} values of $P\circ\chi$. It obtains their equality with the corresponding derivatives of $g$ from equality of the functions on a neighborhood of zero. Thus~\eqref{eq:recurrence} is derived from ordinary differentiable functions; no formal-series coefficients or derivative identities are postulated.

\section{Lean verification and reproduction}
The development uses Lean~4.19.0 and Mathlib revision
\begin{center}
\texttt{c44e0c8ee63ca166450922a373c7409c5d26b00b}
\end{center}
(the official Mathlib \texttt{v4.19.0} release). The complete development is supplied in one source file,\par
\texttt{ThermalKurtosis.lean}. Its main theorem is
\begin{center}
\texttt{ThermalCounterexample.thermal\_kurtosis\_counterexample}.
\end{center}
This theorem has no hypotheses. It assembles the explicit thermal rate conditions with the existence of a $C^4$ dominant spectral branch and the exact negative ratio. The companion theorem \texttt{scgf\_strictly\_dominant} supplies strict dominance over every distinct eigenvalue.

\begin{center}
\begin{tabular}{>{\raggedright\arraybackslash}p{0.38\linewidth}>{\raggedright\arraybackslash}p{0.54\linewidth}}
\toprule
Formal component & Representative theorem names\\
\midrule
Generator and thermal conditions & \texttt{Q\_row\_sum}, \texttt{Q\_offdiag\_pos}, \texttt{hot\_detailed\_balance}, \texttt{cold\_detailed\_balance}\\
Physical heat tilt & \texttt{thermalTilt\_eq\_tilted}\\
Complex spectrum & \texttt{thermal\_eigenvalue\_iff\_cubic}, \texttt{thermal\_dominant\_unique}\\
Smooth branch & \texttt{scgf\_zero}, \texttt{scgf\_smooth}, \texttt{scgf\_equation}\\
Derivative calculation & \texttt{cubic\_comp\_iterated}, \texttt{cumulant\_equations}, \texttt{scgf\_deriv\_four}\\
Negative normalized cumulant & \texttt{kurtosis\_exact}, \texttt{kurtosis\_neg}\\
\bottomrule
\end{tabular}
\end{center}
The source contains no \texttt{sorry}, \texttt{admit}, custom axiom declarations, or native-computation proof shortcuts. The final axiom inspection reports only Lean/Mathlib's standard foundational dependencies: propositional extensionality, classical choice, and quotient soundness. These are distinguished from an assumption of any part of the conjecture or its conclusion. Compiler output is supplied in \texttt{verification.log}. The source hashes, compiler command, and exit status are recorded in \texttt{source-checks.json}.

To reproduce, install Lean through the standard \texttt{elan} toolchain manager, clone Mathlib, check out the pinned revision, and obtain its official compiled dependency cache. The bundle's \texttt{README.md} gives exact commands; \texttt{verify.sh /absolute/path/to/mathlib4} runs the Lean checker and the independent rational checker. The essential proof command, from the pinned Mathlib directory, is
\begin{verbatim}
lake env lean /absolute/path/to/ThermalKurtosis.lean
\end{verbatim}
The proof is checked against the displayed definitions, real numbers, derivatives, and the complex spectrum. No numerical tolerance enters the result.

The optional \texttt{model\_check.py} uses only Python's standard library and exact integer/fraction arithmetic. It expands the characteristic determinant as a Laurent polynomial and also differentiates the directly tilted eigenvector equation, using the normalization $\pi r_n=0$ for $n\geq1$. Both methods agree on all four fractions. This independent arithmetic calculation is a cross-check; Lean supplies the proof of existence, spectral identification, and differentiation.

\section{Relation to previous work and limits}
Three-level thermal models, selective bath couplings, and characteristic-polynomial methods are established. Segal~\cite{segal2018} derives a tilted generator and its cubic characteristic equation for a three-level thermal refrigerator. That work gives explicit mean-current and noise formulas. The model and calculation method used here belong to this established framework.

Negative fourth cumulants of general driven Markov currents are also known. Gerry and Segal~\cite{gerry2023} obtain them in biased modular chains. Their uniform directional bias does not supply the static-energy, bath-resolved thermal realization required by Eq.~(19). The target paper also permits negative fourth cumulants for general unicyclic networks. The result here supplies an explicit thermal model with this sign.

Holubec and Novotn\'y~\cite{holubec2019} study fourth current cumulants in a thermal refrigerator with classical and coherent regimes. Saryal et al.~\cite{saryal2019} study thermal transport and higher cumulants in quantum and non-Markovian settings. Ptaszy\'nski and Esposito~\cite{ptaszynski2025} examine heat fluctuations near phase transitions, including equilibrium fourth cumulants. These results concern related models or regimes; they do not imply nonnegativity for the model in Theorem~1.

Salazar~\cite{salazar2022} studies entropy-production skewness under a detailed fluctuation theorem. Delvenne and Van Brandt~\cite{delvenne2024} derive entropy-production moment bounds under their stated assumptions. These skewness and moment results are distinct from the fourth spectral heat cumulant considered here. An independent-increment assumption cannot be transferred to general multistate current dynamics without proof.

The negative ratio shows that a negative fourth spectral heat cumulant alone cannot exclude the displayed classical thermal model class. The witness has positive $C_3$ and does not refute the separate two-bath skewness conjecture, Eq.~(20). It gives no minimal-state theorem, optimal negative bound, or robustness theorem for added bath channels. The bath couplings are selective, although every total off-diagonal rate is positive. The model is out of equilibrium and does not address a statement restricted to equilibrium.

The result concerns an exact finite model. Formal verification covers the matrix spectrum and derivatives. The identification with long-time stochastic cumulants remains outside the formalization, as specified in Section~1. No claim of empirical validation or priority is made.

\section*{Research and publication}
This work is fully autonomous agentic research. The initial research and proofs were generated autonomously by GPT 6 with Pro Reasoning in ChatGPT. GPT 6 Astra in Codex autonomously produced the formalizations and corrections. The authors checked the proofs to the best of their mathematical abilities. This check is not expert validation or peer review. Dev D. Goyal is responsible for publication.

\clearpage
\begingroup\setlength{\parskip}{3pt}
\begin{thebibliography}{9}\setlength{\itemsep}{4pt}
\bibitem{ptaszynski2022} K. Ptaszy\'nski, ``Bounds on skewness and kurtosis of steady-state currents,'' \emph{Phys. Rev. E} \textbf{106}, 024119 (2022). \href{https://doi.org/10.1103/PhysRevE.106.024119}{doi:10.1103/PhysRevE.106.024119}; \href{https://arxiv.org/abs/2204.08421}{arXiv:2204.08421}, version 4. Target: Eq.~(19); thermal spectral model: Eqs.~(63)--(70).
\bibitem{segal2018} D. Segal, ``Current fluctuations in quantum absorption refrigerators,'' \emph{Phys. Rev. E} \textbf{97}, 052145 (2018). \href{https://doi.org/10.1103/PhysRevE.97.052145}{doi:10.1103/PhysRevE.97.052145}; \href{https://arxiv.org/abs/1803.03713}{arXiv:1803.03713}.
\bibitem{gerry2023} M. Gerry and D. Segal, ``Random walks on modular chains: Detecting structure through statistics,'' \emph{Phys. Rev. E} \textbf{108}, 024135 (2023). \href{https://doi.org/10.1103/PhysRevE.108.024135}{doi:10.1103/PhysRevE.108.024135}; \href{https://arxiv.org/abs/2305.09759}{arXiv:2305.09759}.
\bibitem{holubec2019} V. Holubec and T. Novotn\'y, ``Effects of noise-induced coherence on the fluctuations of current in quantum absorption refrigerators,'' \emph{J. Chem. Phys.} \textbf{151}, 044108 (2019). \href{https://doi.org/10.1063/1.5096275}{doi:10.1063/1.5096275}; \href{https://arxiv.org/abs/1905.10894}{arXiv:1905.10894}.
\bibitem{saryal2019} S. Saryal, H. Friedman, D. Segal, and B. K. Agarwalla, ``Thermodynamic uncertainty relation in thermal transport,'' \emph{Phys. Rev. E} \textbf{100}, 042101 (2019). \href{https://doi.org/10.1103/PhysRevE.100.042101}{doi:10.1103/PhysRevE.100.042101}; \href{https://arxiv.org/abs/1907.10767}{arXiv:1907.10767}.
\bibitem{ptaszynski2025} K. Ptaszy\'nski and M. Esposito, ``Critical heat current fluctuations in Curie-Weiss model in and out of equilibrium,'' \emph{Phys. Rev. E} \textbf{111}, 034125 (2025). \href{https://doi.org/10.1103/PhysRevE.111.034125}{doi:10.1103/PhysRevE.111.034125}; \href{https://arxiv.org/abs/2411.19643}{arXiv:2411.19643}.
\bibitem{salazar2022} D. S. P. Salazar, ``Thermodynamic Skewness Relation From Detailed Fluctuation Theorem,'' \emph{Phys. Rev. E} \textbf{106}, L042101 (2022). \href{https://doi.org/10.1103/PhysRevE.106.L042101}{doi:10.1103/PhysRevE.106.L042101}; \href{https://arxiv.org/abs/2208.11206}{arXiv:2208.11206}.
\bibitem{delvenne2024} J.-C. Delvenne and L. Van Brandt, ``Moments of Entropy Production in Dissipative Devices,'' \emph{Phys. Rev. Lett.} \textbf{134}, 237101 (2025). \href{https://doi.org/10.1103/jqsx-ycjv}{doi:10.1103/jqsx-ycjv}; \href{https://arxiv.org/abs/2404.00422}{arXiv:2404.00422}.
\end{thebibliography}
\endgroup
\end{document}
